\begin{statement}{P3F.41}{Proved}
For each of the four marked sources and the eight source/root sectors in the proof, let $M=\mathcal O(aG-O)$ and $Q=bG$. For every indicated root order $m$ and every nonzero integer $n$ with $\gcd(m,n)=1$, the displayed maps construct a proper holomorphic family with scalar $\lambda=r^m$, one simple whole-fibre $I_1$ zero, and marked cap meridian $mq+n\gamma_{\rm out}$. Its orbit has fundamental group $\mathbb Z/|n|$ and the integral homology displayed in \textup{(SC.4)}. Its total space is the standard oriented $S^6$ exactly for $n=\pm1$. The actual line maps and every gluing degree are retained below.
\end{statement}
\begin{proof}
\par

\medskip\noindent\textbf{SC.1. Explicit sources and the parameter domains}

The following source models retain their actual marked section:

\par\medskip{\small\noindent
\begin{tabularx}{\linewidth}{@{}XXX@{}}
\toprule
Source & Equation and G & Fibres \\
\midrule
E7-A & y\textsuperscript{2}=x\textsuperscript{3}+tx+1, G=(0,1) & III* at infinity, three I1 \\
E7-B & y\textsuperscript{2}=x\textsuperscript{3}+tx-t, G=(1,1) & III* at infinity, II at0, I1 at-27/4 \\
E6-A & y\textsuperscript{2}=x\textsuperscript{3}-3x+t\textsuperscript{2}+2, G=(1,t) & IV* at infinity, I2 at0, I1 at(\pm)2i \\
E6-B & y\textsuperscript{2}=x\textsuperscript{3}+tx+t\textsuperscript{2}, G=(0,t) & IV* at infinity, III at0, I1 at-4/27 \\
\bottomrule
\end{tabularx}
}\par\medskip

The first two are verified in P3F.37/P3K.24. For E6-A the discriminant
is 27t\textsuperscript{2}(t\textsuperscript{2}+4). The factorization
(y-t)(y+t)=(x-1)\textsuperscript{2}(x+2) gives a rational chart
u=(y-t)/(x-1), v=(y+t)/(x-1), with inverse
x=uv-2, t=(x-1)(v-u)/2, y=(x-1)(u+v)/2.
For E6-B its discriminant is t\textsuperscript{3}(4+27t). Put
u=(y-t)/x, v=(y+t)/x; the inverse is
x=uv-(v-u)/2, t=x(v-u)/2, y=x(u+v)/2.
Thus these are rational elliptic surfaces, with the stated minimal
coefficient orders. At infinity their cubic good roots send G to a
nonzero fixed three-torsion point. At the I2 in E6-A, the two local
branch orders of G are one and one; at the III in E6-B the fourth
good root sends G to the nonzero fixed two-torsion point. Thus G has
nonidentity components in E6 and A1, height 1/6, and G.O=0.
The root discriminant and height identity show that G is primitive
and the Mordell--Weil group is torsion free.

Set M=O(aG-O), Q=bG. Use the following actual data, with a single
cap at infinity, no other logarithmic sites, and one chosen I1 as
the simple whole-fibre zero:

\par\medskip{\small\noindent
\begin{tabularx}{\linewidth}{@{}XXXX@{}}
\toprule
Sources & (a,b) & Root order m & Root type \\
\midrule
E7-A,B & (2,1) & m=4l+3, l(\ge)0 & III \\
E7-A,B & (1,2) & m=4l, l(\ge)1 & smooth \\
E6-A,B & (6,1) & m=3l+2, l(\ge)0 & IV \\
E6-A,B & (2,3) & m=3l, l(\ge)1 & smooth \\
\bottomrule
\end{tabularx}
}\par\medskip

For each entry choose an integer n with gcd(m,n)=1. The actual cap
norm will be A\_r\textsuperscript{m=B\_lambda}n, lambda=r\^{}m, with marked meridian
m q+n gamma\_out. In these domains n is nonzero. Every retained
reducible fibre has translation transitive on its components:
at the retained A1 it is odd b, while a is even. The source Hom
line has whole-fibre degree ab/d=1 (d=2 or6). Its degree follows
also directly from the section intersections: for E7 it is1;
for E6 (6,1), (6G).O=(5G).O=2 and G.O=0; for (2,3),
(2G).O=(-G).O=(-3G).O=0. At each component either M has zero
multidegree or Q fixes the component, so the vertical Hom divisor
is a whole fibre. The selected section of that base line supplies
the prescribed one simple I1 zero.

\par

\medskip\noindent\textbf{SC.2. Root equations and exact affine source norms}

For E7 reduced roots use (F37.1) and the equation in K24 for E7-B.
For E6-A at infinity,
\[
 y_0^2=x_0^3-3q^4x_0+q^4+2q^6,\qquad G=(q^2,q^2).
\]
For E6-B,
\[
 y_0^2=x_0^3+q^3x_0+q^4,\qquad G=(0,q^2).
\]
For m=3l+2, k=2l+1, substitute q=z\^{}m,
x0=z\^{}\{2k\}X, y0=z\^{}\{3k\}Y. The respective equations are
\[
\begin{array}{ll}
Y^2=X^3-3z^{4l+4}X+z^2+2z^{6l+6},&
G'=(z^{2l+2},z),\\
Y^2=X^3+z^{l+2}X+z^2,&G'=(0,z).
\end{array}
\]
They have the reduced IV minimal regular fibre. In the chart
X=zX1,Y=zY1 its two nonidentity component labels are distinguished
by Y1=(\pm)1. The actual deck
\[
 D(X,Y,z)=(\zeta_m^{-2k}X,\zeta_m^{-3k}Y,\zeta_mz)
\]
has character zeta\_m\textsuperscript{\{-(3k+1)\}=zeta\_m}\{-2m\}=1 on Y1.
It fixes the three component labels. G' is nonidentity and 6G'
has identity component. Resolving the displayed A2 surface point
retains these labels and the source translation.

For good roots, take q=z\^{}m with k=3m/4 for E7 and k=2m/3 for E6.
The central smooth curve and its deck have order four or three
respectively. In a coordinate w at any marked fixed section,
the minimal good-root chart has weights (r-1,1), r=4 or3.
The good-root limit of Q is O, since b is the component-group
order (two or three). The limit of aG is a nonzero fixed point
of that order. These are the exact marked limits from the source
coordinates, not unmarked smooth elliptic replacements.

There is a unique analytic m-division R of Q' with the selected
central point. At a reduced root multiplication by m is invertible
on its additive identity group and on its component group, since
m is prime to two or three. At a good root use the local inverse
of multiplication by m at O. Deck invariance follows from
uniqueness. The actual source map
\[
                         a_0=t_{-nR}D
\]
satisfies a0\^{}m=t\_\{-nQ'\}. Translation extends to the actual minimal
regular source by the graph/minimality argument used in F37.

\par

\medskip\noindent\textbf{SC.3. The good-root line map as a toric correspondence}

Let r=3 or4 be the minimal good root. Near a marked root section
choose its holomorphic coordinate w=0. Centering at the actual
equivariant section preserves the diagonal weights (r-1,1).
The quotient lattice is
\[
 N'=\mathbb Z^2+\mathbb Z(r-1,1)/r.
\]
The minimal quotient resolution has rays
\[
 v_j=((r-j)/r,j/r),\qquad 0\le j\le r.
\]
Consecutive cones have determinant 1/r in the ordinary lattice,
hence are regular in N'. The following actual common refinements
are regular in the original root lattice:
\[
\begin{array}{ll}
r=3:&(1,0),(2,1),(1,1),(1,2),(0,1),\\
r=4:&(1,0),(3,1),(2,1),(1,1),(1,2),(1,3),(0,1).
\end{array}
\]
Every adjacent determinant is one. Let beta be the resulting root
blowup and pi its toric map to the quotient minimal resolution.
The underlying lattice inclusion gives every chart map by its
dual monomials. Thus pi is an actual proper generically finite
map; no finiteness of its exceptional fibres is assumed.

The resolved strict section is the divisor of the w-axis ray.
Its integral Cartier support function is
\[
 h_w(v)=\max(v_w-(r-1)v_z,0).
\]
Indeed it is w-(r-1)z on the first minimal cone and zero on the
remaining cones. The character w-(r-1)z belongs to the quotient
dual lattice. Pulling that Cartier function to the common
refinement, while beta pulls the original section with valuation
v\_w, gives the exact exceptional correction
\[
 \beta^*(\hbox{root section})-\pi^*(\hbox{resolved section})
 =\sum_{\text{exceptional }v}\min(v_w,(r-1)v_z)E_v.       \tag{SC.1}
\]
For r=3 the coefficients along the listed exceptional rays are
(2,1,1); for r=4 they are (3,2,1,1,1). Subtract the identical
O-chart contribution from the aG-chart contribution. This gives
an exact isomorphism
\[
 \beta^*N=\pi^*M+D_{aG}-D_O,\qquad N=O(aG'-O),           \tag{SC.2}
\]
including each exceptional coefficient.

At r=3 these are the three A2 quotient points of the IV* model.
At r=4 the two marked fixed points have A3 quotients; the other
two points form an order-two orbit and give the remaining A1
arm. No section divisor in (SC.2) passes through that latter
orbit. The minimal resolved quotients are the original IV* and
III* source models: the invariant Weierstrass substitutions give
their birational maps, and their crepant resolution graphs are
the stated affine E6 and E7 fibres, with no vertical (-1)-curve.
The relative minimal-model graph comparison gives the isomorphism.

The correction is equivariant for the actual suspension. At a
fixed root point, Q'\_0=O and deck equivariance give
\[
 w\longmapsto w\,u(w,z)+z^{r-1}h(z^r),\qquad z\longmapsto z,
 \qquad u(0,0)=1.
\]
This preserves each ideal (w,z\^{}j) and (z,w\^{}j), 1\(\le\)j\(\le\)r-1.
The normalized blowup of their invariant product gives exactly
the displayed simultaneous fan refinement. Therefore the actual
source translation lifts, preserves every exceptional correction
in (SC.2), and gives an extending line map on N commuting with D.
For m=r l, pull this whole correspondence through the additional
base change and resolve the fibre product. The same invariant
ideals, now (w,z\^{}\{jl\}), retain the lifted source map and the
pullback divisor identity. This supplies the actual line and B0,
not only their generic-fibre Picard classes.

The natural deck weights of N at the marked nonzero point and O
are respectively zeta\_r\^{}\{r-1\} and zeta\_r\^{}\{-(r-1)\}, up to a common
character. They are the weights of (SC.1) at the section. The
central Picard class is nonzero of order p=3 for E6 and p=2 for
E7. The exact connecting map
\[
 (\mathbb C/\Lambda)^D\longrightarrow\Lambda/(D-1)\Lambda,
 \qquad[v]\longmapsto[(D-1)v]
\]
identifies it with the primitive p-torsion class of the cap circle
line. A common deck character tensors by a flat line from the
free suspension circle and does not remove this torsion.

At a reduced root the source line is already zero on all original
components. Pullback to the resolved fibre product therefore has
degree zero on every contracted curve and descends through the
actual contractions. Its comparison with N has a whole-fibre
divisor, exactly as in F37, using the reduced III or IV
intersection matrix. This is the complementary actual line map.

\par

\medskip\noindent\textbf{SC.4. Commuting line norms, freeness and holomorphic parameter closing}

The full lift argument is the one verified in IV.3. Lift t\_\{-R\}
first by F, correcting its meromorphic chord section's vertical
divisor by a whole-fibre power. Average by the finite deck character
to arrange DFD\^{}\{-1\}=chi F, chi\^{}m=1. Since mR=Q', write the actual
descended map as B0=beta(z)F\^{}m. It commutes with D by (SC.2);
therefore beta is a nonzero function of q. In particular F and D
both commute with B0. The map F\^{}nD has norm ratio
\[
 (F^nD)^m/B_0^n=\chi^{nm(m-1)/2}\beta(q)^{-n}.
\]
An invariant analytic m-th root of this unit gives A0 with
A0\textsuperscript{m=B0}n and {[}A0,B0{]}=1. The character factor is retained also
for the even good-root orders.

Put A\_r=r\^{}nA0, B\_lambda=r\^{}mB0. Choose u,v with nu+mv=1.
Then
\[
 H=A_r^uB_\lambda^v=rH_0,\qquad H^m=B_\lambda,\quad H^n=A_r.
\]
On a smaller compact source tube H0 has bounded norm ratios, so
H is strictly contracting for small \textbar r\textbar. The cap is exactly
N\^{}times/. Its norm shell proves smoothness and properness,
and its strict contraction proves freeness of the delta circle.
For outward h=H\^{}\{-1\}, the single-A cut gives
\[
                         q\mapsto-nh,\qquad\gamma\mapsto mh,
\]
hence precisely m q+n gamma. All source translations and line
characters are retained in this cap and its annular comparison.

The holomorphic ordinary-to-cap translation cochain is solved
once at c0 and then changed by n(c-c0)s\_infty delta/m.
Its positive cut is the actual A\_r. At c-\textgreater c+m its change is
n s\_infty delta, so the inverse closing is multiplication by
(t-tc)\^{}\{-n\}. The positive deck is its inverse. Near the selected
cusp this multiplier and its inverse are B\_lambda iterates
composed with bounded source translations and invertible units.
They extend through its specified Mumford charts. At the other
fillings they are units. This constructs the proper holomorphic
family over the punctured r-disc, with lambda=r\^{}m, including
negative n.~No source root alone is used to choose a parameter
return class.

\par

\medskip\noindent\textbf{SC.5. Integral source and boundary maps}

All maps below are executed with full Smith transformations in
\path{code/part3/iii_star_good_complement.py} and
\path{code/part3/iv_star_complement.py}. Their helpers are retained in
integer\_chain.py. The E7 integral basis is O,F,E7,G as in F37.
For E6 use O,F,C0,\ldots,C5,A1,G. The E6 graph has edges
(0,1),(1,2),(2,3),(3,4),(2,5), highest-root marks
(1,2,3,2,1,2), and G meets C0 and the A1 component.
The full Gram determinant is -1. Translation by G cyclically
permutes the three affine E6 arms and swaps the two A1 components.
Its section columns come from the exact height and component
formula. The resulting matrix is checked to preserve the Gram
form, and the actual contracting source map is its negative b-th
power. No separate unmarked lattice is substituted.

Let V be the source minus its starred fibre tube. That tube has
k=8 or7 component classes and no odd homology. Its inclusion
matrix L is the affine component matrix, retaining the multiplicity
vector. The exact complement chain model has
\[
 C_4=\mathbb Z,\ C_3=\mathbb Z,\ C_2=\mathbb Z^{10},\
 C_1=\mathbb Z^k,\ C_0=\mathbb Z,\qquad
 d_4=1,\ d_2=L^tB.
\]
The source and relative complexes before taking this cone are
free and even, so their degree-one homotopy group is zero.
The displayed translation and component permutation therefore
give the actual functorial mapping-torus cone. The source boundary
has C3=Z,C2=Z\textsuperscript{k,C1=Z}k,C0=Z,d2=L\^{}tBL; its maps to the complement
are 1,L,1 in the corresponding degrees. These are the chain maps
used in the exported calculation, so all boundary inclusions are
retained.

The source complement is simply connected. In E7-A this is the
three-cycle trace argument of F37, and E7-B has the two vanishing
cycles of II meeting once. E6-A has I2+I1+I1. If its primitive
vanishing cycles spanned an index d sublattice, the trace of their
product would be 2 modulo d\textsuperscript{2} (expand the product of rank-one
transvections, including the repeated I2 cycle). The inverse IV*
monodromy has trace-1, so d\textsuperscript{2} divides3, forcing d=1; rank one is
excluded by the same trace. E6-B has a III fibre whose vanishing
cycles already meet once. The source section then kills all base
meridians. Thus the exterior orbit W\_V has pi1=Z gamma.

For the reduced partners, H2(W\_V)=Z, H3(W\_V)=ZF and H4=0.
The boundary H1 and H2 are Z\textsuperscript{2}, with F=p z1, p=2 or3.
Its H2 inclusion is (0,-1) in a primitive basis. The second
column is the actual relative-Wang cycle
p(P\_\{-G\}-O) plus its integral removed-root correction, and is
a unit. The original line need not be that primitive generator:
for E6 its free coefficient is -2. The cap has H2=ZC, with
F specializing to pC, so z1 maps to a unit C. The other column
is an arbitrary retained integer kappa C. The full degree-two
map is unimodular for every kappa.

For the good partners, use the exported boundary basis
(eta,r,s), with orders (p,0,0), and exterior basis (tau,z), with
orders (p,0). The exact H2 inclusions and fibre classes are:
\[
\begin{array}{c|c|c|c}
 & H_2(E)\longrightarrow H_2(W_V)&F\text{ in }H_2(E)&F\text{ in }H_2(W_V)\\
E7,\ b=2&
\begin{pmatrix}-7&-1&3\\0&0&1\end{pmatrix}&(0,-1,0)^t&(1,0)^t\\
E6,\ b=3&
\begin{pmatrix}-2&2&-2\\0&0&-1\end{pmatrix}&(0,1,0)^t&(2,0)^t .
\end{array}                                                     \tag{SC.3}
\]
The torsion coefficient in each first column is a unit modulo p.
The cap has H2=ZC; F maps with degree one to C by actual parallel
transport through the smooth good-root source. Thus its H2 row
is (0,-1,kappa) for E7 and (0,1,kappa) for E6. The free two-by-two
Mayer--Vietoris block is unimodular, and the torsion block is an
isomorphism Z/p-\textgreater Z/p.~This proves the complete H2 isomorphism
without choosing or deleting kappa.

In degree one the good boundary torsion maps primitively to the
cap torsion. The generic elliptic lattice transports identically
to the central good fibre. The two quotient lattices are
Lambda/(D-1)Lambda and Lambda/(D\^{}u-1)Lambda, with u inverse
to n modulo m. Their sublattices agree integrally: for order3
use D+I=-D\textsuperscript{2}, and for order4 an odd power is D or D\^{}\{-1\}.
The multiplying matrices are unimodular. Eliminate these actual
torsion isomorphisms; the remaining free degree-one matrix is
\[
                         D_1=\begin{pmatrix}0&1\\ n&-m\end{pmatrix}.
\]
The reduced partners have that same free matrix without torsion.

In degree three the primitive boundary basis is the oriented
source boundary followed by F gamma. The exterior row is (0,1).
The norm-circle projection has degrees -n and m on those two
classes. Their transverse source cycle specializes to the primitive
invariant full central fibre, both for the smooth root and for
the cyclically permuted reduced components. Hence the cap row
is (-n,m), and D3=D1. There is no extra source H3 class. Degree
four is Z-\textgreater0; degree zero is the incidence column (1,-1)\^{}t.
These specify every gluing degree.

Van Kampen kills the cap elliptic lattice by its actual incoming
generic elliptic loops, since the exterior source complement is
simply connected. In the reduced cases there is no cap elliptic
fundamental group. The remaining relations identify gamma=m h
and impose n h=0. Thus pi1(Y)=Z/\textbar n\textbar. The displayed integral
Mayer--Vietoris maps give
\[
 H_0(Y)=H_5(Y)=\mathbb Z,\quad
 H_1(Y)=H_3(Y)=\mathbb Z/|n|,\quad H_2(Y)=H_4(Y)=0.       \tag{SC.4}
\]
For n=(\pm)1 the orbit is the standard S5. The sole fixed sphere has
the primitive O(-1)\^{}2 Hopf model, and all other delta actions are
free by the actual norm inequality. The Gysin and fixed-tube maps
of K22 give the standard S6 total space. For nonunit n the
same primitive local Hopf meridian kills the principal-circle
kernel and gives pi1(X)=pi1(Y), so these actual total spaces
have nontrivial fundamental group.

This proves the one-cap sphere sectors of the table with their
full local line maps and all integral orbit gluing degrees.
The complete source, line and integral-map proof has passed the
independent audits recorded above. P3K.26 proves the actual return
through its separate signed page and line comparison.
\end{proof}